% ---- Introduction (headline face — converted verbatim from d11_intro_draft.md) ----
\section*{1. Introduction}

\subsection*{Opening problem}

When an embodied agent accumulates experience, where should that
experience live --- folded into the policy's parameters, or kept outside
and retrieved when needed? The question is old and unresolved. The
parametric view is seductive: a fluent agent, like a fluent human,
should not have to look things up mid-task; skill that has been
practised enough becomes automatic, and automaticity feels like
something written into the substrate. The contextual view is
pragmatic: retrieval systems that keep experience external and look it
up have repeatedly matched or beaten parametric memorisation in
language tasks, at a fraction of the training cost. For embodied agents
the stakes are concrete --- a policy that must run at control rate cannot
carry a growing context window, so ``bake it into the weights'' is not
just philosophically appealing but architecturally convenient. Whether
it \emph{works} is an empirical question that, for embodied skill, has --- to
our knowledge --- not been asked under controlled conditions.

\subsection*{Why this test bed answers it}

We ask it with a design that isolates the one variable that matters. A
memory-augmented teacher --- a VLM that retrieves lessons from its own
past attempts --- produces trajectories on a vision-conditioned
single-arm reach task (an Isaac Lab bimanual rig with the right arm
frozen). From those \emph{same} trajectories we distil student policies
whose training targets differ only in whether retrieved-lesson content
is present, and we evaluate them under matched protocols --- including
one that re-attaches inference-time retrieval to a fixed set of student
weights. The critical pair shares identical adapter weights and differs
only in protocol: memory supplied by baking it into those weights
during distillation, versus memory supplied through a retrieval context
at inference. This turns ``parametric vs contextual memory'' from a
slogan into a controlled contrast on one axis.

\subsection*{Pre-registration as a stance}

Because the headline outcome is a null --- and nulls are the most easily
disbelieved result in machine learning --- we pre-registered in two
locked layers. The hypotheses, comparison arms, primary tests, and
their directionality were registered \emph{before any adapter trained}; the
final analysis rules (the pairing-integrity gate and its mechanical
test-selection fallback) were locked \emph{after collection but before any
$p$-value was computed} --- the raw success counts, visible in collect
logs, were frozen out of the test-selection decision by that locked
rule. The pre-registration is not decoration: it is what lets a reader
take the null as evidence rather than as a failure to find an effect we
wanted. Every result below is tagged confirmatory or exploratory
against that registered plan.

\subsection*{Contributions}

\begin{enumerate}
\item \textbf{A design criterion for where experience should live}: marginal
knowledge --- the correction that holds on average across situations ---
transfers into weights, while conditional knowledge --- the correction
that depends on the current situation --- belongs in context, because
weights can only approximate a conditional function with a constant
whereas context conditions on it for free. We name and operationalise
this marginal/conditional distinction and show it predicts the split
we observe. \emph{(This is the transferable takeaway; the results below
are its evidence.)}
\item \textbf{A pre-registered null on parametric memory} (T1): distilling the
memory-augmented teacher's behaviour into student weights beat an
action-only control by $+1$ pp (n.s.), and we withdrew the
memory-in-weights claim per the registered falsification clause ---
scoped to this recipe (single-round SFT + composable-KTO), this
task, this data scale, and this adapter capacity (LoRA rank-16). A
no-training diagnostic (\S4.4b) finds no conditional trace that more
data of the same kind could amplify.
\item \textbf{A reversal}: on a \emph{fixed} set of student weights, attaching
inference-time retrieval added $+34$ pp (identical-weights contrast,
C\_retrieval vs A\_ctrl\_rat, $z=5.15$, exploratory); the pre-registered
protocol contrast (T4, C\_retrieval vs B\_main, each with its own
adapter) agrees at $+23$ pp ($z=3.36$, confirmatory). Retrieval matches
the teacher's own 49\% success rate at ${\sim}50$--$75\times$ lower inference cost.
\item \textbf{A mechanism} (R1 $\Delta$X probe, pre-registered per-arm prediction):
distillation transfers the \emph{marginal} distribution shift (a static
prior all arms adopt regardless of situation) but not the
\emph{conditional} structure (the situation-specific correction that
retrieval supplies), cleanly selecting the registered H\_behavior
branch --- the correction's carrier is the action distribution, not
the memory content.
\end{enumerate}

\tmlrnote{v1.1 consistency flag (integration): the contribution list is
L0a-only; the abstract now carries an L2 cross-task sentence. At merge,
decide whether to add a 5th contribution for the L2 result (must stay
behind the n.s. line: ``a cross-task probe on a harder task --- direction
agrees, effect n.s./sub-MDE, mechanism replicates'') or leave L2 to the
Experiment-2 section only. PI decision; not silently added here.}

\subsection*{One-sentence preview}

The two headline results are one finding: distillation moves the
average, retrieval supplies the situation --- so the recipe is to distil
the competence and externalise the memory.
